\documentclass[11pt]{article}
\usepackage[margin=1in]{geometry}
\usepackage{amsmath}
\usepackage{enumitem}
\usepackage{hyperref}
\usepackage{titlesec}
\usepackage{parskip}
\titleformat{\section}{\large\bfseries}{\thesection}{1em}{}
\titleformat{\subsection}{\normalsize\bfseries}{\thesubsection}{1em}{}

\title{\textbf{SwitchZip: A Cost-Aware Hybrid Classical/Neural Lossless Compressor}\\
\large Final Year BS Project Proposal}
\author{}
\date{}

\begin{document}
\maketitle

\section*{Is it lossless?}
Yes. SwitchZip only decides \emph{which} lossless compressor to use on each
block of data. Both options (classical and neural) are lossless on their
own, and the choice made per block is stored alongside the data, so
decompression always recovers the original file bit-for-bit.

\section{Key Terms}
\begin{description}[leftmargin=1.2cm, style=nextline]
  \item[Lossless compression] Shrinking a file so it takes less space while
  still being able to reconstruct the exact original afterward. (Lossy
  compression, e.g.\ JPEG or MP3, permanently discards some detail.)
  \item[Compression ratio] Original size $\div$ compressed size. Example:
  a 10~MB file compressed to 2.5~MB has a ratio of 4.
  \item[Entropy coding] Assigning shorter codes to more predictable data
  and longer codes to less predictable data. Huffman coding and arithmetic
  coding are both entropy coders.
  \item[Arithmetic coding] Represents a whole message as one number in
  $[0,1)$. At each step the current interval is split according to the
  predicted probability of each possible next symbol, and the piece
  matching the actual symbol becomes the new interval.
  \item[Predictor] Anything that estimates ``how likely is each possible
  next symbol, given what came before.'' Better predictions shrink the
  interval less on the correct guess, which is what makes the output
  smaller.
  \item[Neural predictor] A predictor built from a neural network (e.g.\ an
  LSTM or a small transformer). Usually predicts far better than simple
  statistical models, at the cost of more computation per symbol.
  \item[Classical compressor] A non-neural compressor such as
  \texttt{gzip}, \texttt{zstd}, or PPMd. Cheap and fast, weaker predictor.
  \item[Block] A chunk of the file (e.g.\ 4--64~KB) compressed as one unit,
  so a program can make a separate decision for each chunk.
  \item[Gatekeeper] The classifier this project builds: for each block, it
  predicts whether the neural predictor's gain is worth its extra time.
  \item[Latency] How long compression/decompression takes to run.
\end{description}

\section{Worked Example: Arithmetic Coding}
Suppose the alphabet is \{A, B, C\} with predicted probabilities
$P(A)=0.7$, $P(B)=0.2$, $P(C)=0.1$, and the message to encode is ``A''.

Start with the interval $[0, 1)$. Split it into three pieces sized by
probability:
\[
A \to [0, 0.7), \quad B \to [0.7, 0.9), \quad C \to [0.9, 1.0)
\]
Since the symbol is A, the new interval becomes $[0, 0.7)$ --- a large
range, so it costs few bits to identify a number inside it. If instead the
predictor had assigned A only $P(A)=0.1$, the interval for A would have
been small (e.g.\ $[0, 0.1)$), costing more bits for the same symbol. This
is exactly why a better predictor produces smaller output: it makes the
interval for the symbol that actually occurs as large as possible.

\section{Worked Example: Routing a File}
Take a log file made of two halves: the first half is highly repetitive
boilerplate (timestamps, fixed field names), the second half is
free-form user comments.

\begin{itemize}[nosep]
  \item Block 1--50 (boilerplate): \texttt{gzip} already compresses this
  well because the patterns are short and repeat constantly. A neural
  predictor would barely improve the ratio, so the gatekeeper routes these
  blocks to the classical compressor.
  \item Block 51--100 (free text): predicting the next word needs
  real language understanding. \texttt{gzip} does much worse here, while a
  small neural model captures word-level structure and compresses it
  further. The gatekeeper routes these blocks to the neural compressor.
\end{itemize}
The output file stores: (compressed block 1--50 via gzip) + (compressed
block 51--100 via neural model) + (a short list saying ``blocks 1--50:
classical, blocks 51--100: neural''). Decompression reads that list first
and applies the matching decoder to each block.

\section{Motivation}
Neural predictors combined with arithmetic coding beat classical
compressors by a wide margin on many data types \cite{deletang2024,bellard2019,goyal2018dzip}.
The cost is speed: a neural model has to produce a prediction for every
single symbol, and this speed/ratio tradeoff is identified as the main
obstacle to using neural compression in practice \cite{kipf2026}. There is
also no standard way yet to decide, block by block, whether paying that
cost is worth it \cite{aitdcc2026}.

\section{Project Plan}
\begin{enumerate}[nosep]
  \item \textbf{Baselines.} Run \texttt{gzip}, \texttt{zstd}, and a small
  open-source neural compressor (e.g.\ DZip \cite{goyal2018dzip}) on
  3--4 public datasets: plain text, genomic data, log files.
  \item \textbf{Labels.} For every block, record which method gave the
  best ratio-for-time result. This becomes training data for the
  gatekeeper.
  \item \textbf{Gatekeeper.} Train a small, fast classifier (simple
  block statistics, or a tiny CNN) to predict that label directly from
  the raw block, without running the neural model first.
  \item \textbf{Pipeline.} Block splitting $\to$ gatekeeper $\to$ route
  $\to$ compress $\to$ measure ratio and wall-clock time end to end.
  \item \textbf{Ablations.} Vary block size; measure how often the
  gatekeeper is wrong and what that costs; test on data types not seen
  during training.
\end{enumerate}

\section{Evaluation}
On held-out data, report: compression ratio vs.\ gzip, zstd, and a
fully-neural baseline; total wall-clock time vs.\ the same baselines;
gatekeeper accuracy; and a ratio-vs-time plot showing where SwitchZip
sits relative to ``always classical'' and ``always neural.''

\section*{References}
\begin{thebibliography}{9}

\bibitem{deletang2024}
G.~Del\'etang, A.~Ruoss, P.-A.~Duquenne, E.~Catt, T.~Genewein, C.~Mattern,
J.~Grau-Moya, L.~K.~Wenliang, M.~Aitchison, L.~Orseau, M.~Hutter, and
J.~Veness.
\newblock Language Modeling Is Compression.
\newblock \emph{International Conference on Learning Representations
(ICLR)}, 2024.

\bibitem{bellard2019}
F.~Bellard.
\newblock Lossless Data Compression with Neural Networks.
\newblock Technical report, 2019.

\bibitem{goyal2018dzip}
M.~Goyal, K.~Tatwawadi, S.~Chandak, and I.~Ochoa.
\newblock DZip: Improved General-Purpose Lossless Compression Based on
Novel Neural Network Modeling.
\newblock \emph{arXiv preprint arXiv:1911.03572}, 2019.

\bibitem{kipf2026}
Kipf et al.
\newblock Waiting to Decompress: The Economics of LLM-Based Compression.
\newblock \emph{Conference on Innovative Data Systems Research (CIDR)},
2026.

\bibitem{aitdcc2026}
The 2026 Algorithmic Information Theory Data Compression Challenge.
\newblock Benchmark report, 2026. \url{https://aitdcc.github.io}.

\end{thebibliography}

\end{document}
